% =============================================================================
% Chapter 2 -- Research and Existing Solutions        (owner: Haaniah Ismail)
% Sources: proposal \S{}2, \S{}4, \S{}5, \S{}10, \S{}17 (rewritten, not copied), vendor
% documentation listed in references/references.bib, worksheet \S{}9.
% Dated hands-on trial notes for CodeRabbit / Rovo are NOT yet gathered --
% every claim below that would require a trial is marked [TODO: verify].
% =============================================================================
\chapter{Research and Existing Solutions}
\label{ch:research}

This chapter records what the team studied before committing to a design, what
the comparable products actually do according to their own published
documentation, and where the remaining gap lies. It is deliberately written so
that a reader can check each claim against a source: vendor capabilities are
cited to the vendor's documentation with the date the page was read, research
claims are cited to the original publication, and anything that has not been
verified by us is marked as unverified rather than stated as fact.

% ---------------------------------------------------------------------------
\section{Research Method / Approach}
\label{sec:research-method}

The investigation had four strands, chosen because each answers a different
question and no single strand could answer all of them.

\paragraph{Product documentation review.}
The closest commercial systems were read at the source: their official
documentation, not their marketing pages. For each product we recorded the
capability described, the conditions under which it applies, and the deployment
model. Every capability attributed to a product in this chapter is traceable to
a citation in \texttt{references/references.bib} and was read on 9 October 2026.
Product features change, so the citations carry access dates.

\paragraph{Literature review.}
Three bodies of work were reviewed: requirements traceability, which is the
oldest part of this problem and establishes that linking artifacts by hand does
not survive contact with real projects~\cite{gotel1994}; automated code review,
where learned models now generate and judge review comments at
scale~\cite{li2022codereviewer}; and evaluation corpora for language models on
real repository tasks~\cite{jimenez2024swebench}, which told us what a credible
benchmark looks like and, equally important, what our own evaluation data must
contain that public corpora do not.

\paragraph{Prior final-year projects.}
The proposal records a screening of the FAST School of Computing FYP repository
against this project's claim, to make sure the contribution had not already been
built by an earlier group. The screening found no project addressing cross-tool
evidence reconciliation, with the closest prior work in workflow automation and
requirements traceability. \textbf{[TODO: verify] --- the repository counts quoted
in the proposal's positioning section have not been re-derived for this report;
re-run the screening and record the query before submission.}

\paragraph{Stakeholder consultation.}
A practising senior software engineer was consulted on how his team handles the
mismatch problem today. The three findings from that conversation are reported
in Chapter~1 and shaped two design decisions that appear throughout this report:
precision is worth more than coverage, and a finding must cite its evidence
because the team already ignores most of the alerts it receives.

\paragraph{Trial notes not yet gathered.}
The proposal commits the team to trialling the comparable products directly and
recording what was observed, with dates, instead of relying on documentation.
As of this writing that has not been done: no dated trial notes for CodeRabbit
or Rovo exist in the repository. All statements about those products below are
therefore documentation-derived. \textbf{[TODO: verify --- run the CodeRabbit and
Rovo trials, save dated notes under \texttt{docs/}, and replace or confirm each
row of Table~\ref{tab:existing-solutions} before submission.]}

% ---------------------------------------------------------------------------
\section{Existing Solutions / Products}
\label{sec:existing-solutions}

\subsection{Automated review tools that read the linked issue}

\textbf{CodeRabbit} is an automated pull-request review service. Its Jira
integration pulls issue details into the review context and, in the vendor's
words, assesses whether code changes address the linked issue's acceptance
criteria~\cite{coderabbit-jira}. The linked-issue workflow is explicit: an issue
URL in the pull request description is what makes the check run, and the vendor
separately documents the assessment procedure, in which each objective from the
linked issue is compared against the pull request's changes and returned as
addressed, not addressed, or unclear~\cite{coderabbit-validation}. Two
properties of that design matter to us. It is a hosted service, so the code and
ticket text are processed outside the customer's boundary, and it is scoped to
the repositories and trackers the service is connected to.

\textbf{Atlassian Rovo Dev} provides a comparable check inside the Atlassian
orbit. Its acceptance-criteria check compares the code in a pull request against
requirements listed in linked Jira work items and reports whether each criterion
is met, missing from the code, or needs manual checking~\cite{rovo-ac}. The
published behaviour also states the limits precisely: the work item is found
from the branch name, falling back to commit messages; criteria are read from
specific fields and headings; and the check is unavailable if the work item
description is empty. The feature is hosted and, in practice, ecosystem-locked
to Jira plus Bitbucket or GitHub.

\subsection{Linking and integration primitives}

\textbf{GitHub's native issue--pull-request linking} is the baseline every
developer already has. A pull request can be linked to an issue manually or
through closing keywords, and merging a linked pull request closes the
issue~\cite{gh-pr-issue}. It is reliable, first-class, and entirely
single-vendor: it says nothing about a design file or about the content of the
requirement.

\textbf{Figma's REST API} exposes files, versions and webhooks, which is what
makes a design tool readable as a versioned source at
all~\cite{figma-api}. Having an API is not the same as having intent: the API
returns nodes and version history, not a statement that a given criterion is
satisfied by a given frame.

\textbf{Field synchronisation and automation products} (for example
two-way sync tools and trigger-based automation) move values between systems on
a schedule or on an event. They are useful plumbing, and they are the reason a
team can believe its tools are connected. They do not hold a relationship model,
they do not read code or tests, and they do not express uncertainty.

\subsection{Research systems}

Learned code review is a mature research line: models are trained on large
corpora of real changes and review comments to estimate change quality, generate
review comments and refine code~\cite{li2022codereviewer}. That work
demonstrates that review text can be produced mechanically; it does not address
where the requirement came from, which version of the design the change was
written against, or whether the produced claim can be checked. On the evaluation
side, SWE-bench assembled thousands of real issue--patch tasks from popular
Python repositories and showed that resolving them is genuinely
hard~\cite{jimenez2024swebench}. It is the closest public corpus to what we need,
and it is also the evidence for our constraint: its cases carry an issue and a
patch, not a Figma version, a Jira workflow state, or a labelled verdict about
whether a change satisfies a criterion.

% ---------------------------------------------------------------------------
\section{Comparison of Existing Solutions}
\label{sec:comparison}

Table~\ref{tab:existing-solutions} compares the systems on the four properties
this project is built on, rather than on feature count. The comparison is read
from documentation; it is not the result of a controlled experiment, and no
empirical superiority over a commercial system is claimed anywhere in this
report.

\begin{longtable}{p{0.16\textwidth} p{0.27\textwidth} p{0.24\textwidth} p{0.27\textwidth}}
\caption{Comparison of the reviewed systems against the four properties UNIFY is built on.}\label{tab:existing-solutions}\\
\hline\textbf{System} & \textbf{What it does} & \textbf{Stated limits} &
  \textbf{What UNIFY has to add and prove} \\
\hline\endfirsthead
\hline\textbf{System} & \textbf{What it does} & \textbf{Stated limits} &
  \textbf{What UNIFY has to add and prove} \\
\hline\endhead
CodeRabbit~\cite{coderabbit-jira, coderabbit-validation}
  & Automated PR review that pulls Jira issue context in and returns
    addressed / not addressed / unclear per objective.
  & Hosted; code and ticket text leave the customer boundary; the verdict is
    text in a comment, not a record with resolved citations.
  & Local execution; a packet built deterministically and stored; every citation
    resolved against that packet or the result is downgraded. \\
Rovo Dev~\cite{rovo-ac}
  & Compares PR code with linked Jira work-item requirements; reports met,
    missing or needing manual check.
  & Hosted; Jira-centric; the work item is located from the branch name or
    commit messages; empty description disables the check.
  & Cross-vendor identity that includes the design version; reconciliation of
    late, duplicate and out-of-order events; an auditable case rather than an
    assistant reply. \\
GitHub issue--PR linking~\cite{gh-pr-issue}
  & First-class link between a pull request and an issue, with closing
    keywords and manual linking.
  & Same vendor only; no requirement content, no design, no tests.
  & Relationships that cross tools, each edge carrying how it was established
    and how confident that resolver was. \\
Figma REST API~\cite{figma-api}
  & Programmatic read of files, nodes, version history and webhooks.
  & Returns structure and history, not intent; no notion of a criterion or a
    change under review.
  & Approved design treated as versioned intent linked to a criterion, with
    divergence detection when the implementation moves. \\
Sync and automation tools
  & Move field values between tools on a schedule or a trigger.
  & No relationship model, no code or test evidence, no way to express
    uncertainty about what was copied.
  & Typed edges with provenance, and a link-eligibility rule that stops an
    unconfirmed inference from producing a decisive finding. \\
Learned review models~\cite{li2022codereviewer}
  & Generate review comments, refine code and estimate change quality from a
    diff.
  & No provenance of the requirement; no version identity; output is
    persuasive prose rather than checked citations.
  & A bounded pass budget, a schema-constrained output, and a deterministic gate
    that rejects uncited confidence. \\
Manual review (the baseline)
  & A reviewer opens four interfaces and reasons from memory.
  & Correct when it works; nothing records what evidence supported the
    decision, so the reasoning dies at merge.
  & A durable evidence case per change, re-readable against the artifact
    versions that were current at the time. \\
\hline\end{longtable}

\paragraph{Criteria used for the comparison.}
Five criteria recur in the table and are the ones the evaluation will be judged
against: \emph{scope} (does the system see code, tests, requirements and design
together, or only part of them); \emph{version identity} (can it say \emph{which}
version of each artifact its conclusion came from); \emph{treatment of
uncertainty} (does it abstain, and is the abstention a recorded outcome);
\emph{deployment} (whose infrastructure processes the customer's code); and
\emph{auditability} (can a decision be re-read later, or is it a message that
scrolls away).

% ---------------------------------------------------------------------------
\section{Limitations and Gaps}
\label{sec:gaps}

The overlap with existing products is conceded rather than argued away. Two of
them already perform criterion-level checking against a linked ticket, and one
of them already returns an ``unclear'' state instead of forcing a verdict. It
would be dishonest, and easily disproved, to claim that abstention or
Jira-aware review is novel here.

What none of the reviewed systems does is preserve the evidence behind a
conclusion. The gap this project addresses has four parts:

\begin{enumerate}
  \item \textbf{Versioned provenance across tools.} A conclusion is tied to the
        exact versions of the requirement, the design frame, the diff and the
        tests it was computed from, and those versions are immutable, so the
        conclusion can still be read correctly after the sources move on.
  \item \textbf{Deterministic gating.} A decisive status is produced only when
        every citation resolves against a stored, bounded packet. A fluent
        answer with a citation that does not resolve is downgraded, and the
        downgrade reason is recorded. Confidence is not a substitute for
        evidence.
  \item \textbf{Reconciliation of asynchronous sources.} Requirements, designs
        and code change on three independent clocks with imperfect delivery.
        Duplicate, late and missing events must converge to one final state
        before any comparison is meaningful; the products above assume the
        notification arrived.
  \item \textbf{Local execution.} The reviewed systems are hosted, which
        disqualifies them for organisations whose policy forbids sending
        repository content outside their boundary. Local inference is what makes
        the privacy claim checkable rather than promised.
\end{enumerate}

A fifth, quieter gap is evaluative rather than functional: none of the reviewed
systems publishes a labelled corpus on which its criterion-level claims can be
measured independently. That is why the dataset work in Chapter~8 exists.

% ---------------------------------------------------------------------------
\section{Proposed Solution}
\label{sec:proposed-solution}

UNIFY is a locally deployable change-assurance system built around one
decision: for a single proposed change, does the available evidence support the
intent that change claims to implement? It subscribes to the three source roles
(a design tool, an issue tracker, a change platform), normalises their events
into one envelope, and stores every artifact version as an immutable snapshot.
Relationships between artifacts are persisted as typed edges that record their
own provenance and confidence, so an explicit reference in a pull request title
is distinguishable in the data from a similarity-based inference that no person
has confirmed.

For each atomic acceptance criterion the system assembles a bounded evidence
packet deterministically, runs a fixed-pass local verification loop over it, and
passes the structured result through a gate that resolves every citation against
the stored packet. The output vocabulary is closed: supported, contradicted,
partial, or insufficient evidence, with freshness and review-required carried as
separate fields rather than folded into the verdict. Deterministic rules run
independently of the model and continue to run when the model server is
unavailable.

The architecture, the schema and the process flows are in Chapters~4 and~5; the
part implemented so far in this iteration is in Chapter~6.

% ---------------------------------------------------------------------------
\section{Complex Computing Problem}
\label{sec:ccp}

The complex computing problem is this: how can a system maintain
provenance-preserving relationships among design, requirement, code and test
artifacts that are versioned independently and arrive asynchronously, and then
determine whether one proposed change remains consistent with its linked intent
using only locally executed inference, while keeping false links and unsupported
alerts low when the available evidence is incomplete, delayed, duplicated,
out of order or contradictory?

Five properties make this a research problem rather than an integration task:

\begin{itemize}
  \item \textbf{Identity across heterogeneous tools.} ``The change'' means
        different things in a design file, a ticket and a repository, and each
        tool versions its own record. Resolving them to one unit of work without
        discarding version identity is an entity-resolution problem with no
        shared key.
  \item \textbf{Unreliable asynchronous delivery.} Webhooks are best effort and
        polling is eventually consistent. Correctness has to be defined as
        convergence of the final state under duplicate and out-of-order
        delivery, not as correct handling of a clean stream.
  \item \textbf{Bounded context over large inputs.} A single pull request can
        implicate more code than fits in a model's context. The selection of
        what the verifier may see must be deterministic, budgeted and recorded,
        which makes retrieval part of the correctness argument rather than a
        performance optimisation.
  \item \textbf{Grounding under a hallucination risk.} The component that
        interprets evidence is probabilistic. The system must therefore be
        designed so that an unverifiable claim cannot become a finding: gating,
        not prompting, is the control.
  \item \textbf{Evidence for the evaluation itself.} No public corpus labels a
        design version, a criterion, a diff and a test together. The ground
        truth has to be constructed and independently annotated, which is a
        measurement problem in its own right, with agreement and leakage
        controls as first-class requirements.
\end{itemize}

A routine solution fails on each of these in turn: a synchroniser handles none
of them, a single prompt to a general model handles neither the bounded context
nor the provenance obligation, and a dashboard over the three tools inherits
their delivery guarantees without improving them.

% ---------------------------------------------------------------------------
\section{Technical Challenges}
\label{sec:challenges}

\begin{enumerate}
  \item \textbf{Normalisation without loss of meaning.} One envelope must carry
        three vendors' semantics, including what each considers a state change,
        without silently flattening a distinction that a rule later depends on.
  \item \textbf{Idempotency and reconciliation.} Duplicate delivery must be a
        no-op at the schema level, and a missed event must be synthesised from
        polling rather than left missing. This is unglamorous work and it is a
        precondition for everything above it.
  \item \textbf{Anchor survival.} A requirement pinned to a symbol and line
        range is wrong after the next commit. Re-anchoring by symbol, then
        content hash, then diff-hunk mapping, and marking the anchor broken when
        all three fail, keeps traceability honest instead of approximately
        right.
  \item \textbf{Deterministic packets under a token budget.} Fixed slot budgets,
        a fixed assembly order and a fixed eviction order, with truncation
        recorded in a manifest, so that the same inputs give a byte-identical
        packet and the reviewer can see what was left out.
  \item \textbf{A gate that cannot be talked past.} Schema validation and
        citation resolution must run on the model's output before it is
        displayed, and a verdict citing no evidence span must count as an
        abstention every time, not most of the time.
  \item \textbf{Measuring the system honestly.} Agreement between annotators,
        leakage-free repository-level splits, and baselines that include the
        manual workflow. The team has committed to reporting per-class results,
        including the ones that are unfavourable.
\end{enumerate}

% ---------------------------------------------------------------------------
\section{Research Findings and Design Decisions}
\label{sec:research-findings}

Table~\ref{tab:research-decisions} records the findings that changed the design,
so that a reader can see the chain from source to decision. Each row is
reproducible: the source is cited or is a repository artefact.

\begin{longtable}{p{0.30\textwidth} p{0.30\textwidth} p{0.32\textwidth}}
\caption{Research findings and the design decisions each one produced.}\label{tab:research-decisions}\\
\hline\textbf{Finding} & \textbf{Decision taken} & \textbf{Where it lands} \\
\hline\endfirsthead
\hline\textbf{Finding} & \textbf{Decision taken} & \textbf{Where it lands} \\
\hline\endhead
Commercial tools already check a pull request against a linked
  issue~\cite{coderabbit-jira, rovo-ac}; novelty cannot rest on doing that.
  & Withdraw any ``no product does this'' claim. Compete on local execution,
    versioned provenance, deterministic gating and abstention, each of which is
    measurable.
  & Positioning in this chapter; output vocabulary in Ch.~4; gate in Ch.~5 \\
One of the reviewed tools already returns an ``unclear''
  outcome~\cite{coderabbit-validation}.
  & Do not present abstention itself as novel. Present the mechanism:
    citations resolved against a stored packet, or a recorded downgrade.
  & FR-10; verdict/provenance gate \\
No public corpus labels design version, criterion, diff and test
  together~\cite{jimenez2024swebench}.
  & Build a small, adjudicated oracle ourselves, with blind double annotation,
    agreement statistics and repository-level splits.
  & Dataset package in \texttt{docs/dataset/}; Ch.~8 \\
Stakeholder consultation: more alerts than are read; links tracked
  manually today.
  & Optimise for precision over coverage; abstain rather than guess; route
    confirmation to the code owner instead of the manager.
  & Success criteria in Ch.~1; FR-28; loop split in Ch.~4 \\
Learned review models produce persuasive text without a provenance
  of the requirement~\cite{li2022codereviewer}.
  & Fix the model's job to a bounded, schema-constrained pass over a
    deterministic packet, with a deterministic gate after it; prohibit
    autonomous write.
  & Autonomy ladder, Ch.~5; FR-08, FR-09 \\
Traceability research shows hand-maintained links decay as
  projects evolve~\cite{gotel1994}.
  & Store anchors with commit, symbol and content hash; re-anchor after each
    commit; report anchor survival and mark broken anchors explicitly.
  & FR-22, FR-23; anchor figure in Ch.~5 \\
Figma exposes versions and webhooks but no statement of
  intent~\cite{figma-api}.
  & Treat approved design constraints as structured, versioned evidence linked
    to a criterion, never as pixel comparison.
  & Scope exclusion in Ch.~1; FR-07 \\
\hline\end{longtable}

Two findings did not change the design but bounded it. The documentation review
showed that every comparable product is hosted, which turned local inference
from a preference into a hard constraint that narrows the model choice before
any benchmark is run. And the delivery semantics of the three source roles are
not uniform, which is why polling fallback and reconciliation are architectural
components rather than error handling.
